\documentclass[11pt,a4paper]{article}

\usepackage[utf8]{inputenc}
\usepackage[T1]{fontenc}
\usepackage{lmodern}
\usepackage[margin=2.5cm]{geometry}
\usepackage{amsmath,amssymb}
\usepackage{booktabs}
\usepackage{hyperref}
\usepackage{float}
\usepackage{array}
\usepackage{xcolor}
\usepackage{microtype}
\usepackage{parskip}

\hypersetup{
    colorlinks=true,
    linkcolor=blue,
    citecolor=blue,
    urlcolor=blue
}

\title{\textbf{LindFi: Thermodynamic Finality and\\Physical Identity in Automated Market Makers}}

\author{
    Jorge Pumar\\
    Lindblad Protocol Research\\
    \texttt{founder@lindblad.io}\\
    \texttt{https://lindblad.io}
}

\date{June 2026\\
\small \href{https://doi.org/10.5281/zenodo.20671919}{DOI: 10.5281/zenodo.20671919}}

\begin{document}

\maketitle

\begin{abstract}
We present LindFi, an automated market maker (AMM) operating within the Lindblad Protocol's Spectral Ledger, where the ordering of transactions is anchored in physical entropy rather than validator-controlled block construction. Every liquidity operation --- deposit, withdrawal, and swap --- requires a cryptographic signature generated in real time by physical hardware deriving its key from an SRAM Physically Unclonable Function (PUF). We show that this construction provides structural resistance to three classes of attack that have caused measurable losses in existing DeFi protocols: oracle manipulation, MEV-based front-running, and Sybil-driven liquidity attacks. We further show that the pool state, as an observable of the Spectral Ledger density matrix evolving under the Lindblad master equation, possesses thermodynamic finality --- a completed swap cannot be reversed without violating the second law of thermodynamics. We introduce Physical Coherence Verification (PCV-4) as a yield-weighting mechanism that rewards liquidity providers proportionally to the physical coherence of their hardware, not solely their capital contribution. All results are validated on a live constant-product (x$\cdot$y=k) USDC/USDT pool running on mainnet, including an end-to-end demonstration of a PUF-signed liquidity operation executed by physical hardware.
\end{abstract}

\tableofcontents
\newpage

% ============================================================
\section{Introduction}
% ============================================================

Automated market makers replaced order books with deterministic pricing curves, most commonly the constant-product invariant $x \cdot y = k$ \cite{uniswap}. This design choice solved the liquidity fragmentation problem of order-book exchanges, but it introduced a new attack surface: the pool's state --- and therefore its price --- is a value stored in a smart contract, modifiable by any transaction a validator chooses to include, in whatever order the validator chooses.

This single fact underlies three of the most damaging exploit classes in DeFi history:

\begin{enumerate}
    \item \textbf{Oracle manipulation} --- an attacker temporarily distorts a pool's reserves within a single block to mislead a protocol that reads price from that pool.
    \item \textbf{MEV and front-running} --- a privileged observer of the pending-transaction pool reorders or inserts transactions to extract value from ordinary users.
    \item \textbf{Sybil-driven liquidity} --- liquidity positions carry no identity beyond an address, so capital can appear, extract yield, and disappear without any persistent record of \emph{who} provided it.
\end{enumerate}

LindFi addresses all three by changing what the pool state \emph{is}. In the Lindblad Protocol, the Spectral Ledger --- the layer in which LindFi operates --- is not a smart contract executing on a validator-ordered chain. It is the observable state of a physical system whose evolution is governed by the Lindblad master equation for open quantum systems \cite{lindblad1976}. Every state transition carries a timestamp derived from thermal noise in a Hybrid Stochastic Chua (HSC) circuit, and every transaction is authorized by a signature generated by silicon-derived key material (SRAM PUF) that cannot be extracted or duplicated \cite{pappu2002}.

The result is an AMM where \emph{the question ``what order did these transactions happen in?''} has a physical answer, not a validator-chosen one.

% ============================================================
\section{Background: The Lindblad Cryptography Protocol Stack}
% ============================================================

LindFi inherits its guarantees directly from the four-layer LCP stack described in the base protocol paper \cite{lindblad_paper}:

\begin{itemize}
    \item \textbf{L1 --- SRAM PUF:} Silicon identity. Manufacturing-time variance in SRAM cell power-up states produces a per-chip fingerprint, stabilized via a BCH(255,139,$t$=15) fuzzy extractor (86\% rock-stable bits, 0.00\% intra-device Hamming distance across 12-cycle enrollment).
    \item \textbf{L2 --- P-256 ECDSA:} The fuzzy-extracted key seeds an EVM-compatible signing keypair. The private key is never stored --- it is regenerated from the PUF on each boot.
    \item \textbf{L3 --- Hybrid Stochastic Chua (HSC):} A chaotic circuit driven by thermal ADC noise produces a timestamp that is physically unrepeatable, even by the same hardware.
    \item \textbf{L4 --- Lindblad Consensus:} A dissipative LoRa mesh network propagates state updates; the resulting global state converges to a steady state of the Lindblad master equation.
\end{itemize}

LindFi does not add a fifth layer. It is an application of L1--L4 to a specific state object: the AMM pool.

% ============================================================
\section{LindFi Architecture}
% ============================================================

\subsection{Pool State as an Observable}

A LindFi pool is defined by reserves $(R_a, R_b)$ for tokens $a$ and $b$, an LP-share supply $\Phi$, and a fee parameter $f$. In the Spectral Ledger, these quantities are not stored as plain integers in a database row --- they are expectation values of observables over the system's density matrix $\rho$:

\begin{equation}
    R_a = \text{Tr}(O_{R_a} \cdot \rho), \quad R_b = \text{Tr}(O_{R_b} \cdot \rho), \quad \Phi = \text{Tr}(O_\Phi \cdot \rho)
    \label{eq:observables}
\end{equation}

The density matrix evolves under the Lindblad master equation:

\begin{equation}
    \frac{d\rho}{dt} = -i[H, \rho] + \sum_k \left( L_k \rho L_k^\dagger - \frac{1}{2}\{L_k^\dagger L_k, \rho\} \right)
    \label{eq:lindblad}
\end{equation}

where $H$ is the Hamiltonian governing reversible dynamics and $\{L_k\}$ are Lindblad (jump) operators encoding dissipative, irreversible transitions --- in this context, finalized swaps and liquidity events.

\subsection{The Constant-Product Invariant as a Steady State}

The familiar AMM invariant

\begin{equation}
    R_a \cdot R_b = k
    \label{eq:invariant}
\end{equation}

can be reformulated as a steady-state condition of Eq.~\eqref{eq:lindblad}. Define the pool Hamiltonian $H_{\text{pool}}$ such that its ground-state manifold corresponds to configurations satisfying Eq.~\eqref{eq:invariant} for the current value of $k$. A swap is then not an arbitrary external mutation of $(R_a, R_b)$ --- it is a dissipative transition $\rho \to \rho'$ via a jump operator $L_{\text{swap}}$, after which the system relaxes back onto the steady-state manifold defined by the \emph{new} $k' = R_a' \cdot R_b'$.

This reformulation is what makes Section~\ref{sec:finality} possible: because the transition is dissipative, it inherits the irreversibility of Eq.~\eqref{eq:lindblad}.

\subsection{Operations}

LindFi implements three operations, each a distinct jump operator:

\begin{align}
    L_{\text{add}}&: (R_a, R_b, \Phi) \to (R_a + \Delta a, R_b + \Delta b, \Phi + \Delta\Phi) \\
    L_{\text{remove}}&: (R_a, R_b, \Phi) \to (R_a - \Delta a, R_b - \Delta b, \Phi - \Delta\Phi) \\
    L_{\text{swap}}&: (R_a, R_b) \to (R_a + \Delta_{\text{in}}, R_b - \Delta_{\text{out}})
\end{align}

For $L_{\text{swap}}$ with input token $a$ and fee rate $f$ (in basis points):

\begin{equation}
    \Delta_{\text{out}} = \frac{R_b \cdot \Delta_{\text{in}}(1-f)}{R_a + \Delta_{\text{in}}(1-f)}
    \label{eq:swap}
\end{equation}

The fee component $f \cdot \Delta_{\text{in}}$ is retained in the pool, increasing $k$ and distributing value to LP-share holders proportionally --- modulated by PCV-4 as described in Section~\ref{sec:pcv4}.

% ============================================================
\section{Thermodynamic Finality}
\label{sec:finality}
% ============================================================

\subsection{The Reversibility Problem in Conventional AMMs}

In a validator-ordered blockchain, a swap is finalized when the block containing it is sufficiently deep in the chain that reversing it is \emph{computationally} expensive --- an attacker would need to out-mine or out-stake the rest of the network. This is a probabilistic, economic guarantee. It has a cost, and costs can be paid.

\subsection{Dissipative Irreversibility}

In LindFi, finality is not probabilistic. The jump operators $\{L_k\}$ in Eq.~\eqref{eq:lindblad} describe an open quantum system coupled to an environment --- physically, the thermal bath sampled by the Chua HSC circuit (L3). The dissipative term

\begin{equation}
    \mathcal{D}[\rho] = \sum_k \left( L_k \rho L_k^\dagger - \frac{1}{2}\{L_k^\dagger L_k, \rho\} \right)
    \label{eq:dissipator}
\end{equation}

drives $\rho$ toward a unique steady state $\rho_{ss}$ satisfying $\mathcal{D}[\rho_{ss}] = 0$. Once a swap's jump operator $L_{\text{swap}}$ has acted and the system has relaxed to the new steady state, recovering the pre-swap state $\rho$ from $\rho'$ would require running the dissipative evolution \emph{backward in time} --- decreasing the von Neumann entropy $S(\rho) = -\text{Tr}(\rho \ln \rho)$ of the combined system-environment state.

This is forbidden by the second law of thermodynamics. The cost of reversing a LindFi swap is not "more computation than the attacker can afford" --- it is "less entropy than the universe permits."

\subsection{Implication for MEV}

Maximal Extractable Value strategies --- front-running, sandwich attacks, and reordering --- depend on an adversary's ability to observe a pending transaction and insert their own transaction \emph{before} it, in a chosen order, before either is finalized. In LindFi:

\begin{enumerate}
    \item There is no public mempool of pending Spectral Ledger transactions to observe.
    \item The timestamp determining order is generated by L3 (Chua HSC) at the moment of signing --- a chaotic, thermally-driven value that is unpredictable even to the signing hardware itself before the moment of measurement.
    \item Reordering a transaction after its $L_k$ has acted requires reversing Eq.~\eqref{eq:dissipator}, which Section~4.2 shows to be thermodynamically forbidden.
\end{enumerate}

An adversary wishing to front-run a LindFi swap would need to predict thermal noise in advance --- a physical impossibility, not merely a cryptographic one.

% ============================================================
\section{PUF-Signed Liquidity Positions}
\label{sec:puf}
% ============================================================

\subsection{Authorization Flow}

Every LindFi operation --- \texttt{/pool/add}, \texttt{/pool/remove}, \texttt{/pool/swap} --- requires a completed signature from hardware registered to the requesting LD-address. The flow is:

\begin{enumerate}
    \item The client requests a challenge for a node: \texttt{POST /request-sign \{nodeId, challenge\}} $\to$ \texttt{signId} (status: \texttt{pending})
    \item The physical node polls \texttt{GET /poll-sign?nodeId=...}, detects the pending challenge, and signs it using its L1+L2 key material
    \item The node submits the signature: \texttt{POST /submit-sign \{signId, r, s, puf, chuaNonce\}}
    \item The client polls \texttt{GET /get-sign?signId=...} until \texttt{status: complete}
    \item The LindFi operation is submitted with the completed \texttt{signId}; the fullnode verifies that the signature's reported PUF matches an authorized node for the requesting LD-address via \texttt{device\_authorizations}
\end{enumerate}

\subsection{Live Validation}

This flow was executed end-to-end on mainnet using physical node \texttt{LD95A821D}. A challenge was issued:

\begin{verbatim}
POST /request-sign {"nodeId":"LD95A821D",
                     "challenge":"lindfi_LD-EC6C1139FA3CD393_..."}
-> {"signId": "d68fac19", "status": "pending"}
\end{verbatim}

The physical node detected the pending request, generated a signature using its PUF-derived P-256 key, and submitted:

\begin{verbatim}
r = 61b46777...  (32 bytes)
s = ...           (32 bytes)
puf = 95A821D9AC091073
chuaNonce = ...
\end{verbatim}

The fullnode verified that \texttt{95A821D9AC091073} corresponds to a node authorized for the requesting address, and the operation --- a liquidity addition of 5 USDT for a 4.19 USDC swap --- was executed and recorded in the pool's swap history.

\subsection{Threat Model and Limitations}

The current implementation (Level 1) verifies that a \emph{completed signature with a matching PUF} exists for the requesting address. It does not yet perform full ECDSA $(r,s)$ verification against a derived public key (Level 2) --- this requires standardizing the message encoding and hash function used by the firmware's signing routine, and is planned as a dedicated upgrade. Level 1 already establishes that \emph{some} interaction with the correct physical hardware occurred within the operation's time window; Level 2 will additionally guarantee that the signature is cryptographically unforgeable.

% ============================================================
\section{Physical Coherence Verification (PCV-4)}
\label{sec:pcv4}
% ============================================================

\subsection{Motivation}

In conventional AMMs, fee yield is distributed strictly proportionally to LP-share ownership: $\text{yield}_i \propto \Phi_i$. This treats all capital as equivalent regardless of the reliability of the infrastructure behind it. LindFi introduces a second factor: the \textbf{Physical Coherence Verification score (PCV-4)} of the hardware associated with an LP position.

\subsection{Score Composition}

PCV-4 $\in [0, 1000]$ is computed as a weighted composite of four dimensions:

\begin{equation}
    \text{PCV4} = 1000 \times \left( 0.40\,U + 0.30\,L + 0.15\,S + 0.15\,E \right)
    \label{eq:pcv4}
\end{equation}

where:

\begin{itemize}
    \item $U$ (Uptime) --- the fraction of epochs in a lookback window in which the node submitted a valid attestation: $U = \min(1, \text{participated epochs} / \text{window size})$
    \item $L$ (Latency consistency) --- inverse of the average gap between consecutive participated epochs, $L = \min(1, 1/\overline{\Delta_{\text{epoch}}})$, with $\overline{\Delta_{\text{epoch}}} = 1$ corresponding to perfect consecutive participation
    \item $S$ (Spectral consistency) --- the proportion of a node's completed signatures whose reported PUF matches the node's dominant (registered) PUF identity, $S = n_{\text{matching}} / n_{\text{total}}$, computed over the node's most recent signatures
    \item $E$ (Entropy quality) --- statistical quality of the node's L3 Chua HSC nonces, computed from the bitstream of recent nonces via three lightweight tests: bit balance ($|0.5 - p_1|$), uniqueness ratio, and bit-transition rate, each contributing to $E \in [0,1]$
\end{itemize}

\subsection{Yield Weighting}

For a pool with LP positions $\{(\text{addr}_i, \Phi_i)\}$, the effective weight of position $i$ is:

\begin{equation}
    w_i = \Phi_i \times \frac{\text{PCV4}_i}{1000}
    \label{eq:weight}
\end{equation}

where $\text{PCV4}_i$ is the score of the hardware node authorized for $\text{addr}_i$ via \texttt{device\_authorizations}. Positions not associated with an active mining node receive a baseline $\text{PCV4} = 500$. Fee yield for epoch $t$ is distributed as:

\begin{equation}
    \text{yield}_i = \text{fee}_t \times \frac{w_i}{\sum_j w_j}
    \label{eq:yield}
\end{equation}

\subsection{Live Scores}

PCV-4 was computed for three physical nodes on mainnet using Eq.~\eqref{eq:pcv4} over a 100-epoch lookback window. Following implementation of the $S$ and $E$ dimensions described above, \texttt{LD95A821D}'s score was recomputed using its full signature history:

\begin{table}[H]
\centering
\caption{Live PCV-4 Scores --- June 2026}
\label{tab:pcv4_live}
\begin{tabular}{@{}lccccc@{}}
\toprule
\textbf{Node} & \textbf{Score} & $U$ & $L$ & $S$ & $E$ \\
\midrule
LD95A821D & 989 & 0.990 & 0.990 & 1.000 & 0.977 \\
LD78C57   & 718 & --- & --- & --- & --- \\
LD484C5   & 765 & --- & --- & --- & --- \\
\bottomrule
\end{tabular}
\end{table}

\texttt{LD95A821D}'s $S=1.000$ reflects that all 7 completed signatures in its history reported the PUF registered for that node, with no discrepancies. $E=0.977$ reflects that its 7 Chua HSC nonces (224 bits total) pass bit-balance, uniqueness, and bit-transition tests close to their ideal values. We note explicitly that 7 samples is a small base for the $S$ and $E$ statistics --- both will converge to more stable estimates as the node accumulates additional signed operations through continued LindFi and Physical Oracle use. \texttt{LD78C57} and \texttt{LD484C5} have not yet completed a PUF-signing flow and therefore retain neutral $S=E=0.5$ pending their first signed operations.

These scores reflect real differences in uptime, epoch-participation consistency, and signature-history quality across the three nodes. A liquidity position backed by \texttt{LD95A821D} receives approximately 98.9\% of the maximum possible fee yield per unit of LP-share.

% ============================================================
\section{Resistance to Documented DeFi Exploit Classes}
% ============================================================

\subsection{Oracle Manipulation}

Protocols such as Aave have suffered losses from attacks in which a low-liquidity pool's price was manipulated within a single block to mislead an external oracle consumer \cite{aave_incident}. In LindFi, there is no external oracle: the price \emph{is} $R_b/R_a$, an observable of $\rho$ (Eq.~\ref{eq:observables}). An attacker wishing to manipulate this value must execute $L_{\text{swap}}$ operations, each requiring a fresh PUF-signed authorization (Section~\ref{sec:puf}) and each immediately relaxing the system toward the steady state defined by Eq.~\eqref{eq:invariant} for the new $k$. There is no "single block" in which a transient, reportable distortion can exist.

\subsection{Cascade Liquidations}

Liquidation races --- where bots compete via gas price to be the first to liquidate an undercollateralized position during a market crash --- amplify price declines and have caused significant value destruction during high-volatility events. In LindFi's architecture, transaction ordering is determined by L3 timestamps rather than gas auctions (Section~4.3), removing the economic incentive structure that produces liquidation races. While LindFi's current implementation does not yet include lending (planned for Phase 4), the ordering mechanism generalizes directly to that use case.

\subsection{Sybil-Driven Liquidity}

Because every LP position requires a PUF-signed authorization from registered hardware (Section~\ref{sec:puf}), and because PCV-4 (Section~\ref{sec:pcv4}) accumulates over time as a function of sustained, verifiable uptime, fabricating many high-yield LP positions requires fabricating many physical devices with long, consistent operational histories --- a cost that scales with physical hardware deployment, not with software deployment.

% ============================================================
\section{Live Network State (June 2026)}
% ============================================================

\begin{table}[H]
\centering
\caption{LindFi Pool State --- USDC/USDT, Live Mainnet}
\label{tab:pool_state}
\begin{tabular}{@{}ll@{}}
\toprule
\textbf{Parameter} & \textbf{Value} \\
\midrule
Pool ID & USDC-USDT \\
Reserve $R_a$ (USDC) & 1,014.95 \\
Reserve $R_b$ (USDT) & 1,215.00 \\
LP supply $\Phi$ & 1,099,962,277 (micro-units) \\
Fee rate $f$ & 5 bps (0.05\%) \\
Price $R_b/R_a$ & 1.197107 \\
TVL & \$2,229.95 \\
\bottomrule
\end{tabular}
\end{table}

A complete operational cycle was validated: pool creation, initial liquidity provision, a 100-USDT swap (output: 90.87 USDC), a PUF-signed 5-USDT swap executed via the LindFi web interface with live hardware authorization from \texttt{LD95A821D}, and a subsequent PUF-signed liquidity addition (10 USDC + 10 USDT), all reflected in the live pool state above.

% ============================================================
\section{Future Work}
% ============================================================

\begin{itemize}
    \item \textbf{Level 2 PUF verification:} Full ECDSA $(r,s)$ verification against keys derived from registered node PUFs.
    \item \textbf{PYCO/USDT pool:} Activation pending sufficient organic liquidity; the initial price ratio set by the first liquidity provider is an irreversible economic decision requiring careful timing.
    \item \textbf{PCV-4 sample size:} As nodes accumulate more signed operations, $S$ and $E$ estimates for \texttt{LD78C57} and \texttt{LD484C5} will move from their neutral defaults to measured values, and \texttt{LD95A821D}'s estimates will gain statistical power beyond its current 7-signature history.
    \item \textbf{Lending with PCV-4 collateral:} Extending the L3-ordering argument of Section~6.2 to liquidation queues.
    \item \textbf{LindFi SDK:} Enabling third-party Arbitrum protocols to integrate physically-identified liquidity pools.
\end{itemize}

% ============================================================
\section*{References}
% ============================================================

\begin{thebibliography}{9}

\bibitem{uniswap}
Adams, H., Zinsmeister, N., \& Robinson, D. (2020).
\textit{Uniswap v2 Core.}
Uniswap.

\bibitem{lindblad1976}
Lindblad, G. (1976).
\textit{On the generators of quantum dynamical semigroups.}
Communications in Mathematical Physics, 48(2), 119--130.

\bibitem{pappu2002}
Pappu, R., et al. (2002).
\textit{Physical one-way functions.}
Science, 297(5589), 2026--2030.

\bibitem{aave_incident}
Various (2020--2022).
\textit{Post-mortem analyses of oracle manipulation incidents in DeFi lending protocols.}
Aggregated industry incident reports.

\bibitem{lindblad_paper}
Pumar, J. (2026).
\textit{Lindblad Protocol: Physics-Enforced Hardware Consensus.}
Lindblad Protocol Research.
Zenodo. \href{https://doi.org/10.5281/zenodo.20620319}{DOI: 10.5281/zenodo.20620319}

\bibitem{pyco_paper}
Pumar, J. (2026).
\textit{PYCO: A Physical Coherence Token Emerging from the Lindblad Protocol.}
Lindblad Protocol Research.
Zenodo. \href{https://doi.org/10.5281/zenodo.20643853}{DOI: 10.5281/zenodo.20643853}

\end{thebibliography}

\vspace{2em}
\hrule
\vspace{1em}
\noindent\small\textit{Lindblad Protocol --- The hardware decides. The physics guarantees. The chain records.}

\end{document}
